% P8 limitations
\section{Limitations}
\label{sec:p8-limitations}

\paragraph{Custom synthetic dataset.} Our benchmark is generated by
\texttt{make\_classification} \citep{pedregosa2011sklearn}, not drawn from
production transactions. It reproduces the class imbalance and anonymized
feature shape of single-institution card-fraud data but none of its temporal
drift, verification latency, seasonality, or adversarial adaptation
\citep{dalpozzolo2018realistic}. The current artifact-backed AUCs
($0.7955$, $0.48268$) are therefore results on a stylized testbed and must not
be quoted as production expectations. The margin is dataset- and
environment-specific.

\paragraph{Historical numbers are records, not reproductions.} We could not
reconstruct the environment of the 21~June~2026 internal run that produced the
older $0.975$/$0.948$ pair, so those numbers are historical records, not
headline evidence. The paper's architectural conclusions are
insulated from this by design: the seat assignment rests on the latency, cost,
calibration, and injection arguments of \secref{sec:p8-scorer}, which apply at
accuracy parity and do not depend on the margin or its direction.

\paragraph{No cross-institution validation.} Everything here is one
configuration of one synthetic generator standing in for one institution's
feature view. Fraud patterns, feature pipelines, and base rates vary sharply
across issuers and geographies; the head-to-head has not been replicated on
any second data distribution, public benchmark, or real institution.

\paragraph{Cost estimates are internal.} The $\sim$$85\times$ agentic-triage
cost advantage (\secref{sec:p8-gap-agentic}) is an internal estimate built
from June-2026 token prices, assumed per-alert investigation times, and loaded
analyst compensation. Each assumption is contestable; none has been validated
in a production deployment; the figure is reported to one significant figure
and should be read as ``one to two orders of magnitude,'' not as a measured
constant. The latency and per-transaction cost contrasts of
\secref{sec:p8-scorer} are likewise operational characterizations, not
benchmarked service-level measurements.

\paragraph{Single LLM family, single recipe.} The $0.48268$ challenger is one
instruction-tuned LLM family under one serialization format and one
fine-tuning budget. We did not sweep model
families, scales, serializations, or prompting strategies, and stronger or
differently trained LLMs could narrow the accuracy gap. We note, however, that
the architecture's assignment of seats does not hinge on the gap: the
latency, cost, calibration, and injection arguments of \secref{sec:p8-scorer}
apply at accuracy parity.

\paragraph{Capability gaps argued, not all measured.} Of the four taxonomy
entries, only the scorer head-to-head carries a number of our own. The sensor,
scribe, and cold-start seats are grounded in external literature and primary
regulatory text rather than in our own end-to-end evaluations; the hybrid
architecture of \secref{sec:p8-architecture} is a design contribution whose
integrated performance is future work, not a deployed and audited system.
Nothing in \secref{sec:p8-compliance} is legal advice.
